\documentclass[11pt]{article}
\usepackage[T1]{fontenc}
\usepackage{lmodern,microtype,amsmath,amssymb,amsthm,mathtools,booktabs,array,xcolor,geometry,hyperref,xurl}
\setlength{\emergencystretch}{2em}
\geometry{letterpaper,margin=1in}
\definecolor{ink}{HTML}{18364A}
\hypersetup{colorlinks=true,linkcolor=ink,urlcolor=ink,citecolor=ink}
\newtheorem{theorem}{Theorem}[section]
\newtheorem{lemma}[theorem]{Lemma}
\newtheorem{corollary}[theorem]{Corollary}
\theoremstyle{remark}\newtheorem{remark}[theorem]{Remark}
\DeclareMathOperator{\ord}{ord}
\DeclareMathOperator{\Ind}{Ind}
\DeclareMathOperator{\Reg}{Reg}
\DeclareMathOperator{\Gal}{Gal}
\DeclareMathOperator{\Fix}{Fix}
\DeclareMathOperator{\Ree}{Re}
\newcommand{\Q}{\mathbb Q}
\newcommand{\C}{\mathbb C}
\newcommand{\one}{\mathbf 1}
\newcommand{\Hcal}{\mathcal H}
\newcommand{\eps}{\varepsilon}
\title{\color{ink}\bfseries Beyond Dirichlet GRH:\\Cyclic-generator prime races and linear auxiliary windows}
\author{Research note and reproducible transfer analysis}
\date{25 November 2026}
\begin{document}
\maketitle
\begin{abstract}
We study the compact Alp\"oge--Claude proof of Dirichlet GRH in
\texttt{lsernau/grh-lean}, and refine two of its transferable mechanisms.
First, building on an earlier ProveIt Sylow-filtered reduction, we show
that RH for an arbitrary Galois number field is equivalent, given
Dirichlet GRH, to a one-sided smoothed prime-only estimate supported on
cyclic subgroups generated by chosen generators of Sylow subgroups of
the commutator group. The proof uses Heilbronn characters and monomiality
of finite $p$-groups, not general Artin holomorphy. For symmetric and
alternating groups the test uses only complete splitting and explicitly
listed prime-order factorization patterns. An $S_8$ example has six
patterns with weights $22336,-16,-32,-40,-4,-1$. A quantitative version
bounds the multiplicity of any missed Dedekind zero. Second, an abstract
Gamma-moment argument strengthens the source's quadratic logarithmic
window to a linear one, with a common averaging measure for finitely
many errors and exponentially large coefficients. Both are proved here
using standard inputs. The arithmetic covolume estimate for the new
Frobenius weights remains open; no new unconditional Dedekind RH or
other famous conjecture is claimed. Novelty of the precise refinements
is provisional.
\end{abstract}
\clearpage
\tableofcontents

\section{Scope and source audit}
The upstream repository \cite{upstream} exports five results in
\texttt{grh-compact/src/Main.lean}, with statements in
\texttt{src/Statements.lean} that import only Mathlib. The key endpoint is
\begin{quote}
\small For every positive modulus $q$, Dirichlet character $\chi$ modulo
$q$ and complex $s$ with $\Ree s>1/2$, $L(s,\chi)\ne0$, except for the
principal-character pole at $1$.
\end{quote}
The other endpoints concern the zeta function and the critical line,
including imprimitive Dirichlet characters. This is not a statement
about all Hecke, Artin or automorphic $L$-functions. The public source
snapshot is commit \texttt{72fead102f663d5c7a407964a13dc0cf9d977207}, Lean
4.35.0 and Mathlib \texttt{56a005855ea17e0b0958f73ec801ae140b0a1679}.
The README reports successful kernel/Comparator checks with only
\texttt{propext}, \texttt{Classical.choice}, and \texttt{Quot.sound};
independent build reports \cite{mit,sonoda} substantiate the intended
scope and report checks. The upstream README also cautions that there
has not been a full human reading of the proof. We treat its Dirichlet
endpoint as available, as requested. We inspected selected source
files through the GitHub connector; we did not run a full Lean build
or independently replay the upstream kernel. The present proofs are
ordinary mathematical arguments, not formalized Lean theorems.

\subsection{The transferable architecture}
In the compact proof let $D$ be admissible data, $m$ the number of theta
lifts, $r_D$ the dimension of the row space, $B_D$ the error budget,
$H_D$ the common main height and $R_D$ the remainder. In
\texttt{New/Main.lean}, the proof combines:
\begin{enumerate}
\item Theorems A and B, \texttt{Data.short\_forms} and
\texttt{Data.enough\_forms}, which produce many integral forms of
controlled norms and index, with a near-spanning family of rank $n$;
\item Theorem C, \texttt{Data.integral\_forms\_long}, a lower bound on
its covolume, using fitted twists, height inequalities and the product
formula;
\item The exact remainder identity \texttt{Data.remainder\_eq}, then
\nolinkurl{IsPolar.exists_window_sub_mul_le} and
\texttt{entireL\_ne\_zero\_of\_primeError\_le}.
\end{enumerate}
Schematically, if $V$ is a log squared covolume bound, $\delta$ a log
index bound and $L_D=\log b\,B_D$, then
\[
 r_D(H_D-3L_D)\le V+2\delta,
 \qquad V\le n(H_D-mR_D+L_D).
\]
The missing dimension and index errors are controlled, and the common
main term cancels. The actual source proves
\[
 mR_D\le (\log b)^2 B_D,\qquad B_D=fa(\log b)^3.
\]
Here $m\asymp\sqrt{x}$, and $f,a,b$ have the meanings in the source
guide. Its weighted remainder is an exact expression in
\[
 \lambda_{d,M}(\nu)=\Lambda(\nu)\one_{\nu\equiv1\pmod M}
 (1-(\nu/d)).
\]
The use of non-split primes is built into the arithmetic (not merely an
analytic choice). The cyclotomic action identifies the class $1\bmod M$;
replacing this by a nonabelian conjugacy type is a substantive new
arithmetic problem. The transfer below changes the \emph{target}
statistic, not the meaning of an existing theorem.

\subsection{Relationship with prior reports}
The incoming directory \cite{incoming} contained 16 GRH-named ZIP
archives at inspection. We obtained the central directory and inspected
the manuscript headings/abstracts of three archives using the connector's
base64 reads and ZIP decompression. In particular:
\begin{itemize}
\item \texttt{grh\_sylow\_transfer\_2026-11-25.zip}: a signed prime race
using $\Reg_G-\Ind_{P_p}^G\one$ for one Sylow subgroup of $[G,G]$ per
prime, and a two-pattern reduction for $\mathrm{PSL}_2(7)$. The present
work replaces the Sylow subgroup by cyclic subgroups of its generators,
with a strict support reduction in many symmetric groups. We do not
claim to improve the special $\mathrm{PSL}_2(7)$ result.
\item \texttt{grh\_method\_transfers\_2026-11-25.zip}: a cubic-field
target, a fixed-congruence approximation obstruction, sparse sampling
and a Selmer ramification budget. We use different, nonperiodic prime
weights and a different continuous-window argument.
\item \texttt{grh\_transfer\_analysis\_2026-11-25.zip}: one-sided
Laplace criterion, a populous auxiliary-twist family and explicit
$S_3$ target. Its description also identifies an obstruction for pure
Tate twists. The present criterion needs no twist selection.
\end{itemize}
We could not retrieve complete valid bytes for the large
\nolinkurl{grh_beyond_dirichlet_grh_2026-11-25.zip} (no ZIP end record
in the returned content); one attempt to fetch
\texttt{grh\_lean\_transfers.zip} returned 404. That attempt used underscores rather than the hyphens in the actual
archive name; the correctly named archive was not inspected. The remaining
archives were inventoried by name, not read exhaustively. Consequently the
novelty assessment is relative to the reviewed material and the
literature examined, not a certification of priority. See
\texttt{source\_map.tsv} and \texttt{prior\_reports.md} for provenance.

\section{Cyclic generators as detectors of nonabelian zeros}
Let $E/\Q$ be finite Galois with group $G$, $H=[G,G]$, and for every
prime $p\mid |H|$ choose a Sylow $p$-subgroup $P_p$ of $H$ and a
finite generating list $S_p$ of nonidentity elements. Let $\mathcal T$
contain one representative of every $G$-conjugacy class of cyclic
subgroups $\langle g\rangle$ with $g\in S_p$ for some $p$. Repetitions
could be kept instead. Define the real integer-valued character
\begin{equation}\label{eq:W}
 W=\sum_{C\in\mathcal T} W_C,\qquad
 W_C=\Reg_G-\Ind_C^G\one
     =\Ind_C^G(\Reg_C-\one).
\end{equation}
If $H=1$, the problem is already settled by Dirichlet GRH; the empty
sum is allowed, and statements involving a multiplicity maximum below
are understood only when $H\ne1$. All $W_C$ are sums of monomial
characters induced from nontrivial linear characters of $C$.

\subsection{An order-of-vanishing certificate}
For $\rho$ with $0<\Ree\rho<1$, set
\[
 h_\rho=\sum_{\chi\in\operatorname{Irr}(G)}
    \ord_\rho L(s,\chi,E/\Q)\,\chi,\qquad
 d=h_\rho(1)=\ord_\rho\zeta_E.
\]
Here a pole has negative order. Brauer induction gives meromorphic
continuation, and Artin induction/restriction formalism gives
$\langle h_\rho,\Ind_J^G\psi\rangle=\ord_\rho L(s,\psi,E/E^J)$
for each subgroup $J$ and its characters. These are classical
Heilbronn-character facts \cite{heilbronn}; they do not presume Artin
holomorphy. For every $p$-group $P\le G$, all irreducibles are monomial
(as $p$-groups are supersolvable); their $L$-functions over $E^P$ are
Hecke $L$-functions over intermediate fields and are holomorphic at
$\rho$. This includes the trivial character, whose Dedekind zeta
function has no pole in the open strip. Thus $h_\rho|_P$ is an actual
character. Likewise $h_\rho|_C$ is an actual character for every
cyclic subgroup $C$ (including mixed order). In particular $d\ge0$.
The equality for induced characters follows from Frobenius reciprocity
and the multiplicativity of the Artin $L$-function on virtual
characters, so complex irreducibles create no sign ambiguity in the
integer orders.

\begin{theorem}[Cyclic-generator detection]\label{thm:det}
Assume $\langle h_\rho,\Ind_H^G\one\rangle=0$; in particular this
holds for $1/2<\Ree\rho<1$ under Dirichlet GRH. Let $M$ be the maximum
multiplicity with which a conjugacy class of cyclic subgroups occurs
in the combined generating lists $S_p$. Then
\begin{equation}\label{eq:quant}
 0\le d\le 2M\,\ord_\rho L(s,W,E/\Q).
\end{equation}
If the latter order is zero, all irreducible Artin $L$-functions of
$G$ are holomorphic and nonzero at $\rho$.
\end{theorem}
\begin{proof}
For each $p$, realize $h_\rho|_{P_p}$ on a $d$-dimensional complex
representation $V_p$. It may be made unitary by averaging an inner
product. For $g\in S_p$ the dimension defect is
\[
 D(g)=d-\dim V_p^{\langle g\rangle}
 =\langle h_\rho,W_{\langle g\rangle}\rangle\ge0.
\]
Since $V_p^{P_p}$ is the intersection of the fixed spaces for the
generators, subadditivity of codimension gives
$D_p:=d-\dim V_p^{P_p}\le\sum_{g\in S_p}D(g)$.
For any $u\in H$, write $u=\prod_p u_p$ for its commuting primary
components in $\langle u\rangle$. Each $u_p$ is conjugate in $H$ to
an element of $P_p$, so its fixed-space defect in the actual
representation of $\langle u\rangle$ with character
$h_\rho|_{\langle u\rangle}$ is at most $D_p$. Indeed the defect is
$d- |\langle u_p\rangle|^{-1}\sum_{v\in\langle u_p\rangle} h_\rho(v)$,
and after conjugation $V_p^{P_p}$ lies in that fixed space. The fixed
space of $u$ in the cyclic representation is the intersection of
those of the primary components. The eigenvalues of $u$ are roots
of unity, each with real part at least $-1$, so
\[
 d-\Ree h_\rho(u)\le 2\bigl(d-\dim\Fix(u)\bigr)
 \le 2\sum_p D_p.
\]
Averaging over $u\in H$ gives $d$ on the left, by the hypothesis and
Frobenius reciprocity. Conjugate cyclic subgroups have equal $D(g)$.
Consequently
\[
 d\le 2\sum_p\sum_{g\in S_p}D(g)
 \le 2M\sum_{C\in\mathcal T}\langle h_\rho,W_C\rangle,
\]
which is \eqref{eq:quant}. If the last order is zero, $d=0$. For any
$v\in G$, the actual character $h_\rho|_{\langle v\rangle}$ has
dimension $d=0$ and is identically zero. Thus $h_\rho=0$, and all its
irreducible coefficients vanish. Finally, $\Ind_H^G\one$ is the sum
of the linear characters of $G$, whose $L$-functions over $\Q$ are
Dirichlet functions; the stated Dirichlet GRH implication follows.
\end{proof}

The weight at a nonidentity $u\in G$ is explicitly
\begin{equation}\label{eq:weights}
 W_C(1)=|G|(1-|C|^{-1}),\qquad
 W_C(u)=-\frac{|C_G(u)|\,|C\cap u^G|}{|C|}.
\end{equation}
Thus the positive support is only the identity, and negative support
consists of the conjugacy classes meeting the chosen cyclic subgroups.
The mean over $G$ is zero by reciprocity. The sum in \eqref{eq:W} is
not a virtual difference of unknown analytic status: it is a genuine
sum of induced nontrivial linear characters. In particular its
$L$-function is holomorphic throughout the critical strip. When all
$P_p$ have generators of order $p$, all negative support consists of
\emph{prime-order} elements, rather than all prime-power orders.

\section{The analytic criterion, with prime powers removed}
Write $g_\ell$ for Frobenius at an unramified rational prime $\ell$.
Let $L^*(s,W)$ be the Artin product with all ramified rational primes
omitted. Local factors omitted are holomorphic and nonzero in
$\Ree s>0$ (their reciprocal polynomials have roots only on
$\Ree s=0$), so orders in the strip are unchanged. For $\Ree s>1$ put
\[
 F(s)=-\frac{(L^*)'(s,W)}{L^*(s,W)}
 =\sum_{\ell\text{ unram.}}\sum_{k\ge1}
    W(g_\ell^k)\log\ell\,\ell^{-ks}.
\]
The prime-only transform is
\[
 F_1(s)=F(s)-\sum_{\ell\text{ unram.}}\sum_{k\ge2}
    W(g_\ell^k)\log\ell\,\ell^{-ks}.
\]
The subtracted series is normally convergent in $\Ree s>1/2$:
$|W|$ is bounded and, on a closed half-plane $\Ree s\ge1/2+\eta$,
the sum is dominated by a constant times $\sum_{n\ge2}\log n/n^{1+2\eta}$.
All nontrivial Hecke factors are nonzero on $\Ree s\ge1$ by the
classical zero-free boundary theorem. Hence $F_1(1)$ is finite and
real. For $x\ge1$, define
\begin{equation}\label{eq:E}
 \mathcal E_W(x)=\sum_{\substack{\ell\le x\\\ell\text{ unram.}}}
 W(g_\ell)\log\ell\left(\frac1\ell-\frac1x\right)-F_1(1).
\end{equation}
The centering is essential. A different additive constant would
introduce a pole at $1$. Unlike the source's nonnegative weight,
our character has mean zero, so no $\log x$ term occurs.

For $\Ree s>1$ absolute convergence and elementary integration give
\begin{align}
 \Phi(s)&=\int_1^\infty\mathcal E_W(x)x^{-s}\,dx
 =\frac{F_1(s)/s-F_1(1)}{s-1},\label{eq:Phi}\\
 \operatorname{Res}_{s=\rho}\Phi(s)
 &=-\frac{\ord_\rho L(s,W)}{\rho(\rho-1)}.
 \label{eq:res}
\end{align}
Indeed the integral for an individual $\ell$ is
$\ell^{-s}/[s(s-1)]$. Formula \eqref{eq:Phi} is regular at $1$, and
continues meromorphically in $\Ree s>1/2$. At any real pole
$1/2<\rho<1$, the residue in \eqref{eq:res} is positive.
This is precisely the sign needed for the upper-bound version of
Landau's principle in the source.

\begin{lemma}[One-sided Landau principle in the needed form]\label{lem:landau}
Let $E$ be real and locally integrable on $[1,\infty)$ and of at most
polynomial growth. Suppose its Mellin transform $\Phi(s)=\int_1^\infty
E(x)x^{-s}\,dx$ extends meromorphically to $\Ree s>\beta$ and its real
poles there are simple with nonnegative residues. If
$E(x)\le Kx^{\beta-1}$ eventually for a constant $K>0$, then $\Phi$
has no poles in $\Ree s>\beta$.
\end{lemma}
\begin{proof}
Choose $X\ge1$ after which the bound holds. The nonnegative function
$Kx^{\beta-1}-E(x)$ on $[X,\infty)$ has transform
\[
 G(s)=\frac{KX^{\beta-s}}{s-\beta}-\Phi(s)+\int_1^X E(x)x^{-s}\,dx.
\]
Its real abscissa of convergence $\sigma_c$ is finite above, or
$-\infty$. If $\sigma_c>\beta$, Landau's lemma for a nonnegative
Laplace transform forces a singularity at the real point $\sigma_c$.
For completeness, if it were regular there, Taylor expansion about a
real $s_0>\sigma_c$ sufficiently close to $\sigma_c$ would converge
at a real $s_1<\sigma_c$. The series in powers of $s_0-s_1$ has
nonnegative integrated terms (differentiate the integral at $s_0$).
Monotone convergence would then give convergence at $s_1$, impossible.
The only possible singularity is a simple pole with negative residue
(the negative of a positive real residue of $\Phi$). But $G(s)\ge0$
for real $s>\sigma_c$; a negative residue implies $G(s)\to-\infty$.
Thus $\sigma_c\le\beta$. The nonnegative-function transform is
holomorphic for $\Ree s>\beta$ by absolute convergence on smaller
closed half-planes. The displayed relation gives the result.
\end{proof}

\begin{theorem}[Prime-only cyclic-generator criterion]\label{thm:criterion}
Assume Dirichlet GRH. For a finite Galois extension $E/\Q$ the
following are equivalent:
\begin{enumerate}
\item All nontrivial zeros of $\zeta_E$ have real part $1/2$.
\item For every $\eps>0$, the statistic in \eqref{eq:E} satisfies
$\mathcal E_W(x)\le C_{E,W,\eps}x^{-1/2+\eps}$ for all sufficiently
large real $x$.
\end{enumerate}
Moreover, if the estimate is known for a single $0<\eps<1/2$, no
Dedekind or irreducible Artin zero or pole lies in
$1/2+\eps<\Ree s<1$.
\end{theorem}
\begin{proof}
For sufficiency apply Lemma \ref{lem:landau} with $\beta=1/2+\eps$,
using \eqref{eq:Phi}--\eqref{eq:res}. The trivial estimate
$\mathcal E_W(x)=O(1+(\log x)^2)$ suffices for growth. Since a zero
of $L(s,W)$ would produce a nonzero residue, the bound excludes all
its zeros in the indicated strip. Theorem \ref{thm:det} excludes
Dedekind zeros and all Artin divisors there. Let $\eps\downarrow0$
and use the Dedekind functional equation for the left half of the
critical strip.

For necessity use effective Chebotarev under RH for $\zeta_E$
\cite{lo,gm}. For each fixed conjugacy class $A$, the weighted
unramified-prime count is
$\sum_{\ell\le x,\,g_\ell\in A}\log\ell
=(|A|/|G|)x+O_E(\sqrt{x}\log^2 x)$, by partial summation of the
prime-counting form (or the weighted form and omission of powers).
Summing with weight $W$ cancels the main term. Thus
$A_W(x):=\sum_{\ell\le x} W(g_\ell)\log\ell
=O_{E,W}(\sqrt{x}\log^2 x)$, with ramified primes excluded. Partial
summation gives $F_1(1)=\int_1^\infty A_W(t)t^{-2}\,dt$ (first for
$\Ree s>1$, then by the limit at 1). Hence
$\mathcal E_W(x)=-\int_x^\infty A_W(t)t^{-2}\,dt
=O_{E,W}(x^{-1/2}\log^2 x)$, as required.
\end{proof}

\begin{remark}[Scope of the improvement]
This is a support refinement of a \emph{centered signed monomial
race}. It is not a claim of globally minimal support. In fact the
classical completely-split-prime criterion obtained directly from
$-\zeta_E'/\zeta_E$, with its pole main term removed and prime powers
controlled, can test Dedekind RH using only complete splitting. The
signed mean-zero family here is useful for seeking height comparisons
in which a large common term cancels; it trades additional patterns
for an entire induced-character product, makes the abelian reduction
explicit and provides a multiplicity certificate. Nothing says that
the arithmetic estimate is easier in every field.
\end{remark}

\section{A uniform prime-order test for \texorpdfstring{$S_n$ and $A_n$}{symmetric and alternating groups}}
For $n\ge5$, $[S_n,S_n]=A_n$ and $[A_n,A_n]=A_n$.
Write $p^m1^{n-pm}$ for the cycle type with $m$ disjoint $p$-cycles
and all other points fixed.

\begin{theorem}[Sparse cycle patterns]\label{thm:sym}
For $G=S_n$ or $A_n$, the cyclic subgroups in \eqref{eq:W} may be
chosen with exactly one subgroup of each of the following types:
\begin{align*}
 &p^{p^j}1^{n-p^{j+1}}, &&p\text{ odd prime},\quad
 0\le j<\lfloor\log_p n\rfloor;\\
 &2^{2^j}1^{n-2^{j+1}}, &&1\le j<\lfloor\log_2 n\rfloor.
\end{align*}
The support comprises these types and the identity. In particular no
Frobenius element of composite order is needed. There are
$\sum_{p\le n}\lfloor\log_p n\rfloor-1$ nonidentity patterns.
\end{theorem}
\begin{proof}
A Sylow $p$-subgroup of $S_{p^a}$ is the iterated wreath product
$Q_a=Q_{a-1}^p\rtimes C_p$, with $Q_0=1$: the order is
$p^{1+p+\cdots+p^{a-1}}$, which is the $p$-part of $(p^a)!$.
The outer block rotation is a product of $p^{a-1}$ disjoint $p$-cycles.
Recursing in the blocks gives generators consisting exclusively of
rotations of this form at all levels. The base-$p$ expansion of $n$
allows disjoint blocks of sizes $p^a$; the product of these groups is
Sylow by Legendre's formula. For odd $p$ every generator is even, so
this is a Sylow subgroup of $A_n$ as well.

For $p=2$, the Sylow subgroup of $S_n$ decomposes as $B\rtimes R$,
where $B$ is the elementary abelian group of disjoint bottom pair
swaps, and $R$ is generated by higher block rotations, all even.
This is clear from the recursive construction: an upper rotation
permutes whole pairs without flipping their internal labels. Its
intersection with $A_n$ is $B_{\rm even}\rtimes R$. If the bottom
swaps are $b_1,\ldots,b_m$, then $B_{\rm even}$ is generated by
$b_1b_i$ for $i>1$. These have type $2^2$, and every higher rotation
has type $2^{2^j}$, $j\ge1$. Since $n\ge5$, the type $2^2$ is
available, and all listed levels occur in a largest block.

In $S_n$, cyclic subgroups of prime order with a fixed generator cycle
type are conjugate. In $A_n$, the same is true: if the normalizer in
$S_n$ contains an odd permutation, every conjugation can be adjusted
to an even one. For involutions an individual transposition of one
pair centralizes the generator and is odd. For odd $p$ with repeated
$p$-cycles, swapping two cycles by matching their labels is odd and
centralizes the generator; if there are two fixed points their swap
works as well. In the remaining case (one $p$-cycle and at most one
fixed point), identify the cycle labels with $\mathbb F_p$. A
primitive-root multiplier normalizes the subgroup and is odd, being
a $(p-1)$-cycle on nonzero labels. Therefore one representative per
type suffices also for $A_n$.
\end{proof}

The count equals $\pi(n)-1+O(\sqrt n)$: for exponents $k\ge2$, the
number of primes $p$ with $p^k\le n$ is at most $n^{1/k}$, and the
sum over $k\ge3$ is $O(n^{1/3}\log n)$. In the earlier Sylow-weight
construction the support is all nonidentity prime-power order types
in $A_n$, so the new construction is strictly smaller for many $n$.
The table counts cycle types, not individual split conjugacy classes,
and excludes the identity in both columns.
\begin{center}
\begin{tabular}{rrr}\toprule
$n$ & Cyclic-generator patterns & Full Sylow patterns\\\midrule
8&5&8\\10&6&13\\16&9&32\\20&11&50\\30&15&129\\40&18&271\\\bottomrule
\end{tabular}
\end{center}
For $S_n$ one may compute the weights directly from
\eqref{eq:weights}: the coefficient for a selected $p^m$ type is
$-(p-1)p^{m-1}m!(n-pm)!$. For $A_n$ the weights are exactly half of
these (including the identity); in the split-class case half the
nonidentity powers lie in each class, so the same formula applies
with that factor. The integer weights can be divided by their common
positive gcd without changing the asymptotic criterion.

\subsection{A concrete degree-eight target}
Let $f(T)=T^8-T-1$ and $E$ its splitting field. Exact factorizations
modulo $7$, $19$ and $61$ yield respectively types $(8)$, $(1,7)$,
$(1,2,5)$, with all factors distinct (certificates supplied).
Irreducibility modulo 7 gives transitivity. The 7-cycle gives
primitivity: in any invariant block system the block of its fixed
point must be stable under the 7-cycle, so has size 1 or 8. The fifth
power of the element of type $(1,2,5)$ is a transposition. The graph
of conjugates of this transposition has components forming a block
system. It has an edge and primitivity makes it connected; edge
transpositions of a connected graph generate $S_8$. Therefore
$\Gal(E/\Q)=S_8$.

Let $w(\ell)$ be the following weight, determined by the degrees of
the irreducible factors of $f$ modulo $\ell$ (exclude primes dividing
its discriminant; any extra excluded unramified primes merely change
the centering by finitely many terms of analytic transform).
\begin{center}
\begin{tabular}{lr}\toprule
Factorization pattern & $w(\ell)$\\\midrule
$1^8$&$22336$\\
$2^2 1^4$&$-16$\\
$2^4$&$-32$\\
$3\,1^5$&$-40$\\
$5\,1^3$&$-4$\\
$7\,1$&$-1$\\
All other patterns&$0$\\\bottomrule
\end{tabular}
\end{center}
These are the weights of $W/6$, not a claim that $W/6$ is an integral
character. The underlying $W$ is the genuine character; we divide
the scalar error by 6 only. The exact average is zero: the six class
sizes are $1,210,105,112,1344,5760$, respectively. Their weighted sum
vanishes. Denote by $c$ the value at 1 of the analytically continued
prime-only series $\sum_\ell w(\ell)\log\ell\,\ell^{-s}$.
The following explicit, still unproved, arithmetic estimate would
settle RH for this splitting field (and the off-critical-line
nonvanishing of all its Artin factors):
\begin{equation}\label{eq:s8}
 \sum_{\ell\le x}w(\ell)\log\ell(1/\ell-1/x)-c
 \ \le\ C_\eps x^{-1/2+\eps}\qquad(\text{each }\eps>0).
\end{equation}
The sum uses only the six patterns above, and requires no auxiliary
quadratic zero-density argument. The old full-Sylow filter also uses
$4\,2\,1^2$, $4^2$ and $3^2 1^2$; all three disappear here.

\section{A linear auxiliary-scale selection lemma}
The source lemma \nolinkurl{IsPolar.exists_window_sub_mul_le} in
\nolinkurl{New/Endgame/OneSidedCriterion.lean} selects $u\in[t,t^2]$ with
$E(u)-XE'(u)\le e^{-t}$, uniformly for $0\le X\le e^t$.
We prove a replacement on a linear window $[t,Lt]$ (with a
function-dependent constant $L$), simultaneously for finitely many
errors. This result is independent of RH/GRH.

\begin{theorem}[Common truncated Gamma average]\label{thm:window}
Let $E_1,\ldots,E_q:[0,\infty)\to\mathbb R$ be measurable with
$|E_j(v)|\le B e^{bv}$ for constants $B>0$, $b\ge0$.
Suppose for some $a>b$ the Laplace transforms
$F_j(z)=\int_0^\infty E_j(v)e^{-zv}\,dv$ continue holomorphically
to a neighborhood of the closed disk $|z-a|\le a$.
For every $\gamma>0$ there are $L>1$ and $\kappa>0$ such that, for
all sufficiently large $t$, the probability measure
\[
 d\nu_t(v)=\frac{\one_{[t,Lt]}(v)v^{\lceil\kappa t\rceil}e^{-av}\,dv}
 {\int_t^{Lt} u^{\lceil\kappa t\rceil}e^{-au}\,du}
\]
satisfies $|\int E_j\,d\nu_t|\le e^{-\gamma t}$ for every $j$.
If the $E_j$ are continuous, there exist at most $q+1$ points of
$[t,Lt]$ and nonnegative weights summing to one with the same bound
on all their weighted averages.
\end{theorem}
\begin{proof}
Compactness of the disk and finiteness of the family yield a radius
$R>a$ with all transforms holomorphic on a neighborhood of
$|z-a|\le R$. Let $C$ bound their absolute values there. By the Cauchy
estimate and differentiation under the defining integral at $a$,
\[
 \left|\int_0^\infty v^k e^{-av}E_j(v)\,dv\right|
 =|F_j^{(k)}(a)|\le C k! R^{-k}.
\]
Use the Gamma probability measure $d\mu_k=a^{k+1}v^k e^{-av}dv/k!$.
Its full mean is bounded by $Ca(a/R)^k$. For the low tail, using
$e^{-av}$ as a Markov test on $v<t$, one gets
\[
 \int_{v<t}|E_j(v)|d\mu_k\le B e^{(a+b)t}2^{-(k+1)}.
\]
For the high tail fix $h=(a-b)/2>0$. Integrating the bound
$e^{bv}\one_{v>Lt}\le e^{-hLt}e^{(b+h)v}$ gives
\[
 \int_{v>Lt}|E_j(v)|d\mu_k\le B e^{-hLt}
        \left(\frac a{a-b-h}\right)^{k+1}.
\]
For $k=\lceil\kappa t\rceil$, choose $\kappa$ with both
$\kappa\log(R/a)>\gamma+2$ and
$\kappa\log2>a+b+\gamma+2$. Then choose $L>1$ so large that
$hL>\gamma+2+\kappa\log(a/(a-b-h))$. Each of the three contributions
is $O(e^{-(\gamma+2)t})$. The same tail estimates with $B=1$ in the
majorant $e^{bv}$ show that $\mu_k([t,Lt])\to1$. Normalizing the
restricted measure therefore gives the stated bound. For continuous
errors the expectation vector belongs to the convex hull of the
compact image of the interval in $\mathbb R^q$; Carath\'eodory's theorem
gives the finite combination.
\end{proof}

\begin{corollary}[Uniform moving coefficients]\label{cor:window}
Under the hypotheses, for any $A,\eta>0$ a fixed $L>1$ suffices so
that for all large $t$ and every vector of real coefficients
$|c_j(t)|\le e^{At}$, some $u\in[t,Lt]$ satisfies
\[
 \sum_j c_j(t)E_j(u)\le e^{-\eta t}.
\]
The measure is independent of the coefficient vector; the chosen
point need not be independent of it.
\end{corollary}
\begin{proof}
Apply Theorem \ref{thm:window} with $\gamma>A+\eta+1$. The mean of
the combination is in absolute value at most $qe^{(A-\gamma)t}$,
strictly smaller than $e^{-\eta t}$ eventually. If every point of
the interval were above the claimed bound, its mean would be above
it, contradiction.
\end{proof}

\subsection{How it plugs into the existing Lean proof}
For the source's \texttt{IsPolar}, write $z=s-1$. Its growth is
$O(e^v)$ and the transform agrees with $\int E(v)e^{-zv}dv$ for
$\Ree z>1$. It is holomorphic near $\Ree z\ge0$ except that the
formal definition permits an apparent simple pole at $z=0$; for
\texttt{isPolar\_primeError} the explicit formula removes this pole
(the residue at $s=1$ is zero and the numerator vanishes). For these
prime-error instances, take $a=3$, $b=1$ in the new theorem. This uses
only classical boundary nonvanishing and the local continuation, not
the Dirichlet GRH endpoint. The source's generic \texttt{IsPolar}
interface must be checked for removability at 1 in any other use;
our lemma explicitly requires a holomorphic disk. For
$E(v)=\texttt{primeError}(d,M,e^v)$ and the analogous modulus 1
error, Corollary \ref{cor:window} with $A=\eta=1$ gives exactly the
required inequality for the moving coefficient $X\in[0,e^t]$.
With $t=10^6\log x$, one can select
\[
 x^{10^6}\le x_1\le x^{10^6 L},
\]
instead of the upper bound $\exp((10^6\log x)^2)$. For sufficiently
large $x$, $Lt\le t^2$, so the new admissible window lies inside the
one already allowed in \texttt{eventually\_primeError\_le\_rpow};
the asymptotic main inequality and its exponent are unchanged.
Constants depend on fixed modulus and auxiliary field, not uniformly
on their conductors. The compact proof can still use the old window;
this is a portable sharpening, not a new GRH proof or a better
prime-counting exponent. An important advantage is that a single
averaging measure handles any fixed finite family of compatible
errors, potentially allowing multiple comparison terms.

No simultaneous pointwise smallness of all coordinates follows from
small averages (take $E_1(v)=\sin v$ and $E_2(v)=\cos v$; both satisfy
the disk condition since their Laplace poles are at $\pm i$, but
$E_1(v)^2+E_2(v)^2=1$). This rules out a tempting but false stronger
transfer.

\section{A precise research interface for the missing arithmetic}
The two new statements allow the analytic part of a nonabelian
transfer to be separated from the geometry. Fix an extension $E$,
weight $W$ and any finitely many real errors $E_j(v)$ satisfying the
hypotheses of Theorem \ref{thm:window}. Suppose one can construct
admissible data for each large $x$, with $t=A_0\log x$, and establish
for every $u\in[t,Lt]$ a comparison inequality
\begin{equation}\label{eq:startransfer}
 \mathcal E_W(x)\le C_\eps x^{-1/2+\eps}
       +\sum_{j=1}^q c_j(x) E_j(u),
 \qquad |c_j(x)|\le e^{A_1t},
\end{equation}
with coefficients independent of $u$ and $A_0,A_1>0$ fixed. Here $L$
must be at least the constant furnished by the window theorem for
$\eta$ chosen with $A_0\eta>1/2$. Averaging in $u$ removes the second
term at an exponentially small cost; alternatively choose a good
point by Corollary \ref{cor:window}. Thus \eqref{eq:startransfer}
would imply Theorem \ref{thm:criterion}. This is a sufficient
interface, not an available upstream theorem. It accommodates several
comparison scales at once. In particular $E_j$ may be centered
prime-only errors for fixed signed monomial characters: their
transforms are analytic near $\Ree s\ge1$ by the preceding formulas.

\subsection{Why the arithmetic is not a change of notation}
For a nonabelian extension, the weights depend on Frobenius cycle
type. A fixed congruence class modulo an integer cannot determine
these values: on each compatible abelian Frobenius fiber, Chebotarev
allows every conjugacy class in the fiber. For example in $S_8$, the
identity and a 3-cycle have the same abelianization (even sign) but
weights $22336$ and $-40$. Even enlarging the modulus cannot separate
them: the intersection with a cyclotomic extension is abelian, so
both have the same restriction there. A compatible lift to the
compositum exists for each, and Chebotarev realizes it. This is an
information obstruction to merely retuning $N_0$ or the chosen
cyclotomic tops in \texttt{Data.remainder\_eq}.

A concrete next experiment is to use the permutation modules
$\Q[G/C]$ for the cyclic subgroups in \eqref{eq:W} and seek a norm or
index formula whose local valuation is
$|G|\one_{g=1}-\#\Fix(g\text{ on }G/C)$ at unramified Frobenius $g$.
The equality of this expression with $W_C(g)$ is exact. But these
modules alone do not yield the required short, enough and long
forms. Three obligations have to be established:
\begin{enumerate}
\item A new algebraic integral-form construction with local
divisibility governed by these fixed-point defects, including
ramification and prime-power corrections.
\item A quantitative spanning estimate for the new modules, with a
near-full rank and index loss small enough to preserve the
$1/m$ saving.
\item A lower bound on the same determinant after twists, with the
same archimedean main term and a controlled error, so that
\eqref{eq:startransfer} actually results.
\end{enumerate}
The source's fitted Selmer twists do not establish any of these
for the proposed new Frobenius labels. In particular no arbitrary
nonabelian Hecke character is furnished by a Dirichlet character.
The background module \texttt{HeckeContinuation.lean} proves real
continuation near 1 and limited nonvanishing, not general complex
continuation to a half-plane or its GRH. Classical analytic inputs
used in this note must be formalized or imported separately.

\subsection{Priorities beyond the immediate target}
\begin{center}
\begin{tabular}{p{.26\linewidth}p{.64\linewidth}}\toprule
Problem & What the present transfer supplies, and what it does not\\\midrule
Dedekind RH for nonabelian fields & A proved reduction to a sparse
mean-zero prime race; the arithmetic bound is missing.\\
Artin holomorphy & Nonvanishing and absence of poles off the line
would follow for each certified extension; nothing here excludes
poles at critical-line zeros.\\
Artin primitive-root conjecture & Kummer splitting fields in Hooley's
argument need nonabelian estimates, uniform across fields. A
fixed-field criterion, even if proved arithmetically for each field,
does not provide the uniformity required by Hooley's method.\\
General Hecke or automorphic GRH & The window lemma transfers given
analytic boundary regularity; cyclic-generator character arguments
apply to finite Galois representations, not arbitrary automorphic
representations.\\
BSD, Goldbach, Vandiver & No proved new reduction to these problems
here. Twists and positivity alone do not imply a rank formula,
additive correlations, or Selmer vanishing.\\\bottomrule
\end{tabular}
\end{center}

\section{Reproducibility and formalization plan}
\texttt{checks.py} uses exact permutation closure, integer arithmetic,
rational class averages, and SymPy finite-field polynomial arithmetic.
It verifies the stated Sylow generators in small degrees, tabulates
the selected and full-Sylow supports through $n=40$, and checks the
$S_8$ weights and factorizations. Irreducibility of every modular
factor is independently checked by the finite-field Frobenius/gcd
criterion. It does not numerically test or prove any asymptotic.
\texttt{certificate.json} records 70 Sylow cases, 69 with full group
closure enumerated (one above the cap); 36 rows of pattern counts.
The exact test is reproducible with Python 3 and SymPy. The full
mathematical proofs above do not depend on this enumeration.

A Lean port can proceed in layers (see \texttt{formalization.md}):
first finite group fixed-space identities and the Sylow generator
construction; then the complex Artin/Hecke orders and Heilbronn
restriction identities; then the prime-only transform, the
one-sided Landau principle, and the Gamma-window theorem. The source
has some finite-character infrastructure and its Landau lemma,
but the analytic interfaces need careful extension. No Lean code
is claimed to compile. The arithmetic construction of
\eqref{eq:startransfer} is a research problem rather than an omitted
proof term.

\section{Conclusion}
There is a meaningful extension of the methodology, not just a list
of consequences of RH. The cyclic-generator reduction shrinks the
nonabelian signal to prime-order cycle patterns for all symmetric
and alternating groups and supplies a bound on the multiplicity of
any undetected Dedekind zero. The Gamma argument improves auxiliary
scale selection without assuming any Riemann hypothesis. These are
complete intermediate theorems with explicit hypotheses. The most
important remaining step is to construct a height/covolume comparison
for a genuinely nonabelian Frobenius signal; the source as read
contains no such general result. We have not resolved a previously
open famous conjecture, and do not certify global novelty of the
intermediate statements.

\begin{thebibliography}{9}\small\setlength{\itemsep}{0pt}\setlength{\parsep}{0pt}
\bibitem{upstream} L. Sernau et al., \emph{GRH for Dirichlet
L-functions, in Lean 4}, 25 Nov. 2026,
\url{https://github.com/lsernau/grh-lean/tree/72fead102f663d5c7a407964a13dc0cf9d977207}.
\bibitem{mit} A. V. Sutherland, \emph{Shorter GRH Lean proof:
verification files}, \url{https://math.mit.edu/~drew/shortgrh/}.
\bibitem{sonoda} Independent build report,
\url{https://shosonoda.github.io/lean-rh-alpoge/grh.html}.
\bibitem{incoming} ProveIt, incoming report archives,
\url{https://github.com/VladimirReshetnikov/ProveIt/tree/main/docs/incoming}.
\bibitem{heilbronn} R. Foote, M. R. Murty, \emph{Zeros and poles of
Artin L-series}, Math. Proc. Cambridge Philos. Soc. 105 (1989), 5--11;
for survey context see \emph{Heilbronn characters}, Bull. Amer. Math.
Soc., \url{https://doi.org/10.1090/BULL/1492}.
\bibitem{lo} J. C. Lagarias, A. M. Odlyzko, \emph{Effective versions
of the Chebotarev density theorem}, in Algebraic Number Fields,
Academic Press, 1977, pp. 409--464.
\bibitem{gm} L. Greni\'e, G. Molteni, \emph{An explicit Chebotarev
density theorem under GRH}, J. Number Theory 200 (2019), 441--485,
\url{https://doi.org/10.1016/j.jnt.2018.12.005}.
\bibitem{mv} H. L. Montgomery, R. C. Vaughan, \emph{Multiplicative
Number Theory I: Classical Theory}, Cambridge, 2007, \S15.1--15.2.
\bibitem{hooley} C. Hooley, \emph{On Artin's conjecture}, J. Reine
Angew. Math. 225 (1967), 209--220.
\end{thebibliography}
\end{document}
